% Tabla: Análisis de impacto por nivel de rendimiento (tier)
% Experimento 06 - K-Fold CV (5×3 = 15 folds)
\begin{table}[htbp]
\centering
\caption{Impacto del aumento por nivel de rendimiento línea base}
\label{tab:tier_impact}
\begin{tabular}{lcccc}
\toprule
\textbf{Nivel} & \textbf{Rango F1} & \textbf{N Clases} & \textbf{$\Delta$ F1 (pp)} & \textbf{Desv. Est. (pp)} \\
\midrule
\rowcolor{green!20}
\textbf{Bajo} & $<$ 0.20 & 9 & \textbf{+1.21} & 2.76 \\
Intermedio & 0.20 -- 0.45 & 6 & +0.01 & 0.12 \\
Alto & $\geq$ 0.45 & 1 & +0.02 & 0.00 \\
\bottomrule
\end{tabular}
\vspace{0.5em}\footnotesize

\textbullet\ \textbf{Hallazgo clave}: el aumento con LLM beneficia principalmente a clases de nivel bajo.

\textbullet\ Clases destacadas en nivel bajo: ISTJ (+4.5 pp), ENTJ (+4.3 pp), ISFJ (+3.5 pp).

\textbullet\ Clases de alto rendimiento muestran mejoras marginales; el sistema ya captura sus patrones.

\end{table}
